% Regression: t2_eq50_pull — imported current-behavior fixture from the handoff corpus.
% Formula: (schematic):
% T2 — RMP arXiv:2011.12127, fig. fig3-pepe_eq50redarrows (eq. 50):
% pulling a symmetry-twisted virtual string through a G-injective PEPS
% plaquette.
% Formula (schematic):
%   T . [lambda-ring plaquette with the red string leaving NORTH-WEST]
%   = [lambda-ring plaquette with the red string leaving SOUTH-EAST] . T,
% i.e. the extra tensor T carrying the red (symmetry) string can be pulled
% from the top of the plaquette to the bottom; lambda is the on-ring
% weight, barbs on black legs give V (out) vs V* (in).
% Free tier (genuinely non-grid: a diamond of tensors threaded by a closed
% ring plus curved string arcs).  lambda = on-wire-matrix skin (`ring`);
% the ring = straight + hv/vh joins; the string arcs = route=curve joins.
%
% Missing primitive (tenkz-free.code.tex, \tnjoin style resolution): joins
% stroke `bond`/`fused bond` only -- no operator-hued join (no
% /tenkz/join/operator key), so every wire the source draws RED (the ring
% segments' string continuation, the sweeping arcs, the stubs) renders
% black.  The hue is load-bearing: it separates the symmetry string from
% the black lattice legs.
\documentclass[varwidth=19cm, border=6pt]{standalone}
\usepackage{amsmath, amssymb}
\usepackage{tenkz}
\begin{document}
\[
  % LHS: string leaves the plaquette north-west through T
  \begin{tenkzfree}[tensor style=box]
    \tnput[box]{n}{(0.3,0.95)}{}
    \tnput[box]{ei}{(0.95,0)}{}
    \tnput[box]{s}{(0,-0.95)}{}
    \tnput[box]{wi}{(-0.95,0)}{}
    \tnput[ring]{lam}{(-0.5,0.95)}{\lambda}
    \tnput[box]{w}{(-2.1,0)}{}
    \tnput[box]{t}{(0.3,1.9)}{}
    % black lattice legs (barb out = V, barb in = V*)
    \tnjoin{-2.75,0}{w.west}
    \tnjoin[dir]{w.east}{wi.west}
    \tnjoin{wi.east}{-0.5,0}
    \tnjoin{ei.west}{0.5,0}
    \tnjoin[dir]{ei.east}{1.7,0}
    \tnjoin[dir]{t.south}{n.north}
    \tnjoin[dir]{n.south}{0.3,0.4}
    \tnjoin[dir]{s.south}{0,-1.6}
    \tnjoin{t.north}{0.3,2.45}
    % the virtual ring through the plaquette (source: RED), lambda on it
    \tnjoin{lam.east}{n.west}
    \tnjoin[route=hv]{n.east}{ei.north}
    \tnjoin[route=vh]{ei.south}{s.east}
    \tnjoin[route=hv]{s.west}{wi.south}
    \tnjoin[route=vh]{wi.north}{lam.west}
    % the string: T's east stub, the sweep down to W, W's south stub
    \tnjoin{t.east}{1.1,1.9}
    \tnjoin[route=curve, out=180, in=90]{t.west}{w.north}
    \tnjoin{w.south}{-2.1,-0.75}
  \end{tenkzfree}
  \;=\;
  % RHS: string pulled through, leaving south-east through B
  \begin{tenkzfree}[tensor style=box]
    \tnput[box]{n}{(0.3,0.95)}{}
    \tnput[box]{ei}{(0.95,0)}{}
    \tnput[box]{s}{(0,-0.95)}{}
    \tnput[box]{wi}{(-0.95,0)}{}
    \tnput[ring]{lam}{(-0.5,0.95)}{\lambda}
    \tnput[box]{e}{(2.1,0)}{}
    \tnput[box]{b}{(0,-1.9)}{}
    % black lattice legs
    \tnjoin[dir]{-1.7,0}{wi.west}
    \tnjoin{wi.east}{-0.5,0}
    \tnjoin{ei.west}{0.5,0}
    \tnjoin[dir]{ei.east}{e.west}
    \tnjoin{e.east}{2.75,0}
    \tnjoin[dir]{0.3,1.6}{n.north}
    \tnjoin[dir]{n.south}{0.3,0.4}
    \tnjoin[dir]{s.south}{b.north}
    \tnjoin{b.south}{0,-2.45}
    % the virtual ring (source: RED)
    \tnjoin{lam.east}{n.west}
    \tnjoin[route=hv]{n.east}{ei.north}
    \tnjoin[route=vh]{ei.south}{s.east}
    \tnjoin[route=hv]{s.west}{wi.south}
    \tnjoin[route=vh]{wi.north}{lam.west}
    % the string: B's west stub, the sweep up to E, E's north stub
    \tnjoin{b.west}{-0.8,-1.9}
    \tnjoin[route=curve, out=0, in=270]{b.east}{e.south}
    \tnjoin{e.north}{2.1,0.75}
  \end{tenkzfree}
\]
\end{document}
